\documentclass[10pt,a4paper]{amsart}
\usepackage{amsmath,amssymb,amsthm,graphicx}
% Retained preamble for the Stacks Project chapter snapshot in this directory.
% The chapter's own preamble.tex is a copy of the shared Stacks Project preamble; the block
% below is sliced MECHANICALLY, line for line, from it. Nothing was added, removed or
% redefined; the wrapper class, the xr-hyper/hyperref wiring and the externaldocument list
% are omitted because the assembler supplies the document class.
%
% lines taken from preamble.tex: 3, 5, 9, 11, 12, 156, 157, 158, ... (51 lines total)
%
% package coverage restored from the chapter's own preamble (2026-10-13): xy, xy-2cell

\usepackage[all]{xy}
\xyoption{2cell}
\UseAllTwocells

\newenvironment{reference}{\comment}{\endcomment}
\newenvironment{slogan}{\comment}{\endcomment}
\newenvironment{history}{\comment}{\endcomment}
\theoremstyle{plain}
\newtheorem{theorem}[subsection]{Theorem}
\newtheorem{proposition}[subsection]{Proposition}
\newtheorem{lemma}[subsection]{Lemma}
\theoremstyle{definition}
\newtheorem{definition}[subsection]{Definition}
\newtheorem{example}[subsection]{Example}
\newtheorem{exercise}[subsection]{Exercise}
\newtheorem{situation}[subsection]{Situation}
\theoremstyle{remark}
\newtheorem{remark}[subsection]{Remark}
\newtheorem{remarks}[subsection]{Remarks}
\numberwithin{equation}{subsection}
\def\lim{\mathop{\mathrm{lim}}\nolimits}
\def\colim{\mathop{\mathrm{colim}}\nolimits}
\def\Spec{\mathop{\mathrm{Spec}}}
\def\Hom{\mathop{\mathrm{Hom}}\nolimits}
\def\Ext{\mathop{\mathrm{Ext}}\nolimits}
\def\SheafHom{\mathop{\mathcal{H}\!\mathit{om}}\nolimits}
\def\SheafExt{\mathop{\mathcal{E}\!\mathit{xt}}\nolimits}
\def\Sch{\mathit{Sch}}
\def\Mor{\mathop{\mathrm{Mor}}\nolimits}
\def\Ob{\mathop{\mathrm{Ob}}\nolimits}
\def\Sh{\mathop{\mathit{Sh}}\nolimits}
\def\NL{\mathop{N\!L}\nolimits}
\def\CH{\mathop{\mathrm{CH}}\nolimits}
\def\proetale{{pro\text{-}\acute{e}tale}}
\def\etale{{\acute{e}tale}}
\def\QCoh{\mathit{QCoh}}
\def\Ker{\mathop{\mathrm{Ker}}}
\def\Im{\mathop{\mathrm{Im}}}
\def\Coker{\mathop{\mathrm{Coker}}}
\def\Coim{\mathop{\mathrm{Coim}}}
\DeclareMathSymbol{\boxtimes}{\mathbin}{AMSa}{"02}
\def\QCohstack{\mathcal{QC}\!\mathit{oh}}
\def\Cohstack{\mathcal{C}\!\mathit{oh}}
\def\Spacesstack{\mathcal{S}\!\mathit{paces}}
\def\Quotfunctor{\mathrm{Quot}}
\def\Hilbfunctor{\mathrm{Hilb}}
\def\Curvesstack{\mathcal{C}\!\mathit{urves}}
\def\Polarizedstack{\mathcal{P}\!\mathit{olarized}}
\def\Complexesstack{\mathcal{C}\!\mathit{omplexes}}
\def\Pic{\mathop{\mathrm{Pic}}\nolimits}
\def\Picardstack{\mathcal{P}\!\mathit{ic}}
\def\Picardfunctor{\mathrm{Pic}}
\def\Deformationcategory{\mathcal{D}\!\mathit{ef}}
\usepackage{cleveref}
\crefname{originalsection}{original section}{original sections}
\Crefname{originalsection}{Original section}{Original sections}
\makeatletter
\newlabel{sites-cohomology-example-category-to-point}{{08PF}{source}{}{sourcepointer-1}{}}
\newlabel{sites-cohomology-example-category-to-point@cref}{{[originalsection][0][]08PF}{[1][0][]source}{}{}{}}
\makeatother
\crefname{originallemma}{original lemma}{original lemmas}
\Crefname{originallemma}{Original lemma}{Original lemmas}
\makeatletter
\newlabel{sites-cohomology-lemma-compute-by-cosimplicial-resolution}{{08Q9}{source}{}{sourcepointer-2}{}}
\newlabel{sites-cohomology-lemma-compute-by-cosimplicial-resolution@cref}{{[originallemma][0][]08Q9}{[1][0][]source}{}{}{}}
\makeatother
\crefname{originallemma}{original lemma}{original lemmas}
\Crefname{originallemma}{Original lemma}{Original lemmas}
\makeatletter
\newlabel{sites-cohomology-lemma-eilenberg-zilber}{{08QC}{source}{}{sourcepointer-3}{}}
\newlabel{sites-cohomology-lemma-eilenberg-zilber@cref}{{[originallemma][0][]08QC}{[1][0][]source}{}{}{}}
\makeatother
\title{Vanishing of the homology of symmetric powers below the symmetric degree}
\author{The Stacks Project Authors}
\date{stacks-project:cotangent.tex:local-snapshot-20261012; GNU FDL 1.2 or later}
\begin{document}
\maketitle
\setcounter{subsection}{1}
\noindent\textit{Editorial note.} All delivered mathematical text is the author's own, copied from the pinned Stacks Project chapter file cotangent.tex; no statement or proof step was rewritten, reordered or removed.\par\medskip
\noindent\textit{Editorial note.} The statement refers to the notation and assumptions fixed in a cross-chapter example and lemma; both are declared as navigation pointers and the proof works with that notation.\par\medskip
\noindent\textit{Editorial note.} The three references of the counted spans are recorded as non-load-bearing navigation pointers; the argument that consumes them stays inside the counted proof span.\par\medskip
\noindent\textit{Editorial note.} Preamble: the retained preamble of this chapter (sources/stacks-cotangent/retained-preamble.tex, lines 1-63), containing the Xy-pic setup, the comment environments, the theorem declarations and the macros including \textbackslash{}NL.\par\medskip
\noindent\textit{Allowed standard prerequisites (editorial).}\begin{itemize}
\item The cosimplicial resolution computation of homology on the punctual topos. Cohomology on Sites, Example sites-cohomology-example-category-to-point and Lemma sites-cohomology-lemma-compute-by-cosimplicial-resolution, used at proof lines 2283-2287. Fixes the setting of the statement.
\item The Eilenberg-Zilber type formula for lower shriek of a tensor product. Cohomology on Sites, Lemma sites-cohomology-lemma-eilenberg-zilber, used at proof lines 2315-2321. Turns the E\_1 terms into derived tensor products of derived lower shrieks.
\item The Koszul complex resolving the symmetric powers. used at proof lines 2298-2307. Produces the exact complex and hence the spectral sequence.
\item Flatness of F and the vanishing H\_0(C, F) = 0. used at proof lines 2288-2289 and 2325-2329. Give the base cases and the vanishing for wedge powers.
\end{itemize}\medskip
\ifdefined\hypertarget\hypertarget{sourcepointer-1}{}\fi
\noindent\textit{Source reference.} Original section 08PF: focused-seeds/AGP/sources/stacks-snapshot-20261011/tags.txt, source lines 11023--11023. Navigation pointer into the Stacks Project; the standard fact it supplies is applied inside the counted proof span and the argument resting on it stays there.\par\medskip
\ifdefined\hypertarget\hypertarget{sourcepointer-2}{}\fi
\noindent\textit{Source reference.} Original lemma 08Q9: focused-seeds/AGP/sources/stacks-snapshot-20261011/tags.txt, source lines 11052--11052. Navigation pointer into the Stacks Project; the standard fact it supplies is applied inside the counted proof span and the argument resting on it stays there.\par\medskip
\ifdefined\hypertarget\hypertarget{sourcepointer-3}{}\fi
\noindent\textit{Source reference.} Original lemma 08QC: focused-seeds/AGP/sources/stacks-snapshot-20261011/tags.txt, source lines 11055--11055. Navigation pointer into the Stacks Project; the standard fact it supplies is applied inside the counted proof span and the argument resting on it stays there.\par\medskip
\begin{lemma}
\label{lemma-vanishing-symmetric-powers}
Notation and assumptions as in
Cohomology on Sites, Example \ref{sites-cohomology-example-category-to-point}.
Assume $\mathcal{C}$ has a cosimplicial object as in
Cohomology on Sites, Lemma
\ref{sites-cohomology-lemma-compute-by-cosimplicial-resolution}.
Let $\mathcal{F}$ be a flat $\underline{B}$-module such that
$H_0(\mathcal{C}, \mathcal{F}) = 0$.
Then $H_l(\mathcal{C}, \text{Sym}_{\underline{B}}^k(\mathcal{F})) = 0$
for $l < k$.
\end{lemma}

\begin{proof}
We drop the subscript ${}_{\underline{B}}$ from tensor products, wedge powers,
and symmetric powers. We will prove the lemma by induction on $k$.
The cases $k = 0, 1$ follow from the assumptions. If $k > 1$ consider
the exact complex
$$
\ldots \to
\wedge^2\mathcal{F} \otimes \text{Sym}^{k - 2}\mathcal{F} \to
\mathcal{F} \otimes \text{Sym}^{k - 1}\mathcal{F} \to
\text{Sym}^k\mathcal{F} \to 0
$$
with differentials as in the Koszul complex. If we think of this as a
resolution of $\text{Sym}^k\mathcal{F}$, then this gives a first quadrant
spectral sequence
$$
E_1^{p, q} =
H_p(\mathcal{C},
\wedge^{q + 1}\mathcal{F} \otimes \text{Sym}^{k - q - 1}\mathcal{F})
\Rightarrow
H_{p + q}(\mathcal{C}, \text{Sym}^k(\mathcal{F}))
$$
By Cohomology on Sites, Lemma \ref{sites-cohomology-lemma-eilenberg-zilber}
we have
$$
L\pi_!(\wedge^{q + 1}\mathcal{F} \otimes \text{Sym}^{k - q - 1}\mathcal{F}) =
L\pi_!(\wedge^{q + 1}\mathcal{F}) \otimes_B^\mathbf{L}
L\pi_!(\text{Sym}^{k - q - 1}\mathcal{F}))
$$
It follows (from the construction of derived tensor products) that
the induction hypothesis combined with the vanishing
of $H_0(\mathcal{C}, \wedge^{q + 1}(\mathcal{F})) = 0$ will prove what we want.
This is true because $\wedge^{q + 1}(\mathcal{F})$ is a quotient
of $\mathcal{F}^{\otimes q + 1}$ and
$H_0(\mathcal{C}, \mathcal{F}^{\otimes q + 1})$
is a quotient of $H_0(\mathcal{C}, \mathcal{F})^{\otimes q + 1}$
which is zero.
\end{proof}

\end{document}
